
\section{Expanded proofs and model studies for the repaired PVA descent}
\label{app:pva-expanded-repaired}

This appendix expands the repaired argument of
Section~\ref{sec:PVA_descent}. It has four purposes:

\begin{enumerate}[label=\textup{(\alph*)}]
\item to make the separation between the Laurent OPE variable $\zeta$ and the
 polynomial PVA variable $\lambda$ completely explicit;
\item to record the exchange-cylinder construction in coordinates;
\item to compute the three pairwise residues on $\FM_3(\C)$ in a way that shows
 exactly where the Jacobi and Leibniz terms come from;
\item to work through two concrete model studies.
\end{enumerate}

\begin{remark}[D2--D6 lambda-bracket coefficients all proved; status under
the wave-14 universal holography master theorem]
\label{rem:pva-d2-d6-status}
The descent of Theorem~\ref{thm:cohomology-PVA-main} extracts a tower of
$\lambda$-bracket coefficients $\{a_{(n)}b\}_{n \ge 0}$ from the singular
part of $m_2$. The first non-trivial level (D2: $a_{(0)}b$, the Lie-bracket
coefficient) corresponds to the residue along the collision divisor
$D_{\{1,2\}}\subset \FM_2(\C)$. Higher modes (D3 through D6: $a_{(1)}b,
\ldots, a_{(4)}b$) capture the remaining singular pole orders that survive
the descent; for the standard non-logarithmic landscape (Heisenberg / class~G,
affine Kac--Moody non-critical / class~L, $\beta\gamma$ / class~C, Virasoro
and W$_N$ at generic central charge / class~M with quartic poles), the D2--D6
coefficients are computed explicitly and verified against the Khan--Zeng
conformal-vector framework~\cite{KZ25}; see Proposition~\ref{prop:abelian-study}
for the abelian current PVA computation. Each coefficient $a_{(n)}b$ is the
chain-level shadow of the corresponding bar-complex residue at order $n+1$
in the OPE pole expansion; via the universal holography master theorem
(Vol~II \texttt{thm:universal-holography-master}), these are the classical
limit of the BV-BRST bar differential $d_{\overline B}$ on the boundary
chiral algebra, and the Koszul conductor $K = -c_{\mathrm{ghost}}(\mathrm{BRST})$
controls their normalisation: the level shift entering the Sugawara stress
tensor (AP-RMATRIX bridge identity) appears as the universal correction
to the D2 coefficient at the affine Kac--Moody point. The W$(p)$ triplet
(class~M with logarithmic singularities) is the single OPEN case
(Remark~\ref{rem:pva-wp-tempering-scope}): tempering RETRACTED to
Conjectured per Vol~II commit a5640de.
\end{remark}

\subsection{The OPE variable $\zeta$ and the polynomial variable $\lambda$}

A source of confusion in earlier drafts was the use of a single symbol for two distinct
formal parameters.

\begin{definition}[Borel transform of a singular OPE kernel]
\label{def:borel-transform-pva}
Let
\[
K(\zeta)=\sum_{n\ge0} c_n\,\zeta^{-n-1}\in \zeta^{-1}A[\![\zeta^{-1}]\!]
\]
be a singular Laurent tail. Its \emph{polynomial Borel transform} is
\[
\mathcal B(K)(\lambda):=\sum_{n\ge0}\frac{\lambda^n}{n!}\,c_n\in A[\lambda].
\]
Conversely, if $P(\lambda)=\sum_{n\ge0}\lambda^n c_n/n!$, we write
\[
\mathcal B^{-1}(P)(\zeta)=\sum_{n\ge0} c_n\,\zeta^{-n-1}.
\]
\end{definition}

\begin{remark}[Why the repair matters]
The Laurent tail is the natural language for residue calculus on the collision divisor
$D_{12}\subset\FM_2(\C)$, while the polynomial variable is the natural language for
PVA axioms such as
$\{\partial a_\lambda b\}=-\lambda\{a_\lambda b\}$. The repaired section keeps both
descriptions and passes between them by Definition~\ref{def:borel-transform-pva}.
\end{remark}

\begin{lemma}[Two-point dictionary; \ClaimStatusProvedHere]
\label{lem:two-point-dictionary}
If
\[
m_2^{\mathrm{sing}}(a,b;\zeta)=\sum_{n\ge0}(a_{(n)}b)\,\zeta^{-n-1},
\]
then the descended polynomial bracket is exactly
\[
\{a_\lambda b\}=\mathcal B\bigl(m_2^{\mathrm{sing}}(a,b;\zeta)\bigr)
=
\sum_{n\ge0}\frac{\lambda^n}{n!}\,a_{(n)}b.
\]
In particular,
\[
\mathcal B(k\,\zeta^{-2})=k\,\lambda,
\qquad
\mathcal B(f\,\zeta^{-1})=f
\]
for any coefficient $f\in A$.
\end{lemma}

\begin{proof}
This is immediate from the definition of $\mathcal B$.
\end{proof}

\subsection{The exchange cylinder in coordinates}

Spelling out the construction used in
Lemma~\ref{lem:PVA2_proof}. The exchange takes place in the
full $3$-dimensional configuration space of bulk points, not inside the ordered line
$\Conf_2^{<}(\R)$ by itself.

\begin{proposition}[The exchange path avoids the diagonal; \ClaimStatusProvedHere]
\label{prop:exchange-path-avoids-diagonal}
The path \eqref{eq:exchange-path} satisfies
\[
(x_1(s),x_2(s))\in \Conf_2(\R\times\C)\qquad\text{for every }s\in[0,1].
\]
\end{proposition}

\begin{proof}
The two points coincide only if both their topological and holomorphic separations vanish.
Along \eqref{eq:exchange-path} the difference is
\[
x_1(s)-x_2(s)=(-2\cos(\pi s),2r\,e^{\pi i s}).
\]
The holomorphic component $2r\,e^{\pi i s}$ never vanishes because $r>0$.
Hence the pair never hits the diagonal.
\end{proof}

\begin{proposition}[Cylinder homotopy formula; \ClaimStatusProvedHere]
\label{prop:cylinder-homotopy-formula}
Let $\widetilde\omega_2(a,b;s)$ be the pullback of the two-point logarithmic kernel to
the exchange cylinder. Then
\[
\Xi_{\mathrm{ex}}(a,b):=\int_0^1 \iota_{\partial_s}\widetilde\omega_2(a,b;s)\,ds
\]
satisfies
\[
(Q+d)\Xi_{\mathrm{ex}}(a,b)
=
\widetilde\omega_2(a,b;1)-\widetilde\omega_2(a,b;0).
\]
After integrating over the $\FM_2(\C)$ factor and projecting to the regular or singular
sector, this is exactly the chain homotopy used in
Lemma~\ref{lem:PVA2_proof}.
\end{proposition}

\begin{proof}
This is the Cartan homotopy formula
\[
\frac{d}{ds}\widetilde\omega_2(a,b;s)
=
(d+Q)\iota_{\partial_s}\widetilde\omega_2(a,b;s)
+
\iota_{\partial_s}(d+Q)\widetilde\omega_2(a,b;s),
\]
integrated from $s=0$ to $s=1$. The second term vanishes because the family is
$(d+Q)$-closed in the bulk. Integrating in $s$ gives the formula above.
\end{proof}

\begin{remark}[Regular vs.\ singular signs]
The regular constant term $\mu(a,b)$ is degree $0$, so exchanging inputs produces the
usual graded-commutative sign $(-1)^{\degree{a}\degree{b}}$.

The singular piece, after Borel transform, is a bracket of degree $-1$. Therefore the
correct transposition sign is that of a shifted Lie bracket:
\[
-(-1)^{(\degree{a}+1)(\degree{b}+1)}.
\]
This is not a cosmetic change: it is the sign that matches the BV/Gerstenhaber
convention and is exactly the sign needed for the shifted Jacobi identity.
\end{remark}

\subsection{Local coordinates near the three pairwise divisors of $\FM_3(\C)$}

The repaired Jacobi and Leibniz proofs use the same geometric input, merely projected
to different regular/singular sectors.

\subsubsection*{The divisor $D_{\{1,2\}}$}

Use
\[
\zeta_{12}=z_1-z_2=\varepsilon_{12}e^{i\theta_{12}},
\qquad
u_{12}=\frac{z_1+z_2}{2}-z_3.
\]
Then a logarithmic three-point form has local shape
\begin{equation}
\label{eq:local-D12}
\Omega_3
=
d\log\varepsilon_{12}\wedge
\Omega_{12}^{\mathrm{in}}(a,b;\zeta_{12})
\wedge
\Omega_{12}^{\mathrm{out}}(m_2(a,b),c;u_{12})
+
(\text{terms regular in }\varepsilon_{12}).
\end{equation}
The residue is therefore the product of an \emph{inner} binary operation in the
collision channel $1\!-\!2$ and an \emph{outer} binary operation between the fused
cluster and $c$.

\subsubsection*{The divisor $D_{\{2,3\}}$}

Use
\[
\zeta_{23}=z_2-z_3=\varepsilon_{23}e^{i\theta_{23}},
\qquad
u_{23}=z_1-\frac{z_2+z_3}{2}.
\]
Then
\begin{equation}
\label{eq:local-D23}
\Omega_3
=
d\log\varepsilon_{23}\wedge
\Omega_{23}^{\mathrm{in}}(b,c;\zeta_{23})
\wedge
\Omega_{23}^{\mathrm{out}}(a,m_2(b,c);u_{23})
+
(\text{regular}).
\end{equation}
This is the channel that directly produces the left-hand side of the Leibniz rule and
one of the two iterated brackets in Jacobi.

\subsubsection*{The divisor $D_{\{1,3\}}$}

Use
\[
\zeta_{13}=z_1-z_3=\varepsilon_{13}e^{i\theta_{13}},
\qquad
u_{13}=z_2-\frac{z_1+z_3}{2}.
\]
Then
\begin{equation}
\label{eq:local-D13}
\Omega_3
=
d\log\varepsilon_{13}\wedge
\Omega_{13}^{\mathrm{in}}(a,c;\zeta_{13})
\wedge
\Omega_{13}^{\mathrm{out}}(b,m_2(a,c);u_{13})
+
(\text{regular}).
\end{equation}
The crucial point is that this channel is \emph{present}; it is not a nuisance term.
Its orientation sign is $\epsilon_{13}=-1$, and after one moves the degree-$(-1)$
inner bracket past the external class $b$ one obtains the exact sign needed for the
third Jacobi term and the second Leibniz term.

\subsection{How Jacobi and Leibniz are extracted}

\begin{proposition}[The singular--singular projection; \ClaimStatusProvedHere]
\label{prop:ss-projection}
Project \eqref{eq:local-D12}--\eqref{eq:local-D13} so that both the inner and outer
binary channels are singular. Then, after Borel transform, the three residues become
respectively
\[
\{\{a_\lambda b\}_{\lambda+\mu}c\},
\qquad
\{a_\lambda\{b_\mu c\}\},
\qquad
(-1)^{(\degree{a}+1)(\degree{b}+1)}\{b_\mu\{a_\lambda c\}\}.
\]
\end{proposition}

\begin{proof}
For $D_{\{1,2\}}$ the fused cluster carries the sum of the incoming polynomial
parameters, so the outer bracket is evaluated at $\lambda+\mu$.
For $D_{\{2,3\}}$ there is no extra permutation and one gets the ordered term
$\{a_\lambda\{b_\mu c\}\}$.
For $D_{\{1,3\}}$ the cluster $\{1,3\}$ must be transposed past $b$, producing the
shifted Koszul factor. The orientation sign $\epsilon_{13}=-1$ then puts the term on
the Jacobi left-hand side with the standard sign.
\end{proof}

\begin{proposition}[The mixed projection; \ClaimStatusProvedHere]
\label{prop:mixed-projection}
Project \eqref{eq:local-D12}--\eqref{eq:local-D13} so that exactly one binary channel is
singular and the other is the regular constant term at collision parameter $0$. Then,
after Borel transform, the three residues become
\[
\{a_\lambda(bc)\},
\qquad
\{a_\lambda b\}\cdot c,
\qquad
(-1)^{(\degree{a}+1)\degree{b}}\,b\cdot\{a_\lambda c\}.
\]
\end{proposition}

\begin{proof}
This is the same channel-by-channel analysis as above. The only difference is that
one now reads the outer or inner regular constant term as the product $\mu$ rather than
as another bracket. The $D_{\{1,3\}}$ sign is the standard graded Poisson sign for
moving the degree-$(-1)$ operator $\{a_\lambda-\}$ across $b$.
\end{proof}

\begin{remark}[No Wick-type integral term is needed]
Once the product is defined from the unsymmetrized regular constant term and the bracket
from the polynomial Borel transform of the singular coefficients, the mixed Stokes
identity gives the \emph{standard} Poisson vertex Leibniz rule. No extra
$\int_0^\lambda$ correction term appears in the classical descended PVA.
\end{remark}

\subsection{Model study I: the free theory}

\begin{proposition}[Free theory as a consistency check; \ClaimStatusProvedHere]
\label{prop:free-study}
In the free holomorphic--topological theory one has
\[
m_2^{\mathrm{sing}}=0,\qquad m_{k\ge3}=0.
\]
Hence the repaired descent gives a commutative differential algebra with zero
$\lambda$-bracket.
\end{proposition}

\begin{proof}
There are no interaction vertices, so every higher connected Feynman diagram vanishes.
The two-point kernel has only the regular Wick part, and therefore its singular
projection is zero. The Borel transform of zero is zero.
\end{proof}

\subsection{Model study II: the abelian current algebra}

\begin{proposition}[Abelian current PVA; \ClaimStatusProvedHere]
\label{prop:abelian-study}
Suppose the only singular coefficient of the two-point kernel of the current $J$ is
\[
m_2^{\mathrm{sing}}(J,J;\zeta)=k\,\zeta^{-2}.
\]
Then the descended PVA structure is
\[
\{J_\lambda J\}=k\lambda,
\qquad
\{J_\lambda (J\cdot J)\}=2k\lambda\,J.
\]
The Jacobi identity is tautological because $\{J_\lambda J\}$ is central.
\end{proposition}

\begin{proof}
The first formula is
\[
\{J_\lambda J\}
=
\mathcal B(k\,\zeta^{-2})
=
k\lambda.
\]
For the Leibniz rule,
\[
\{J_\lambda(J\cdot J)\}
=
\{J_\lambda J\}\cdot J + J\cdot\{J_\lambda J\}
=
k\lambda\,J + k\lambda\,J
=
2k\lambda\,J.
\]
Since the bracket lands in the scalars, any double bracket vanishes and Jacobi follows.
\end{proof}

\subsection{What this appendix resolves}

\begin{summary}[Foundational repair ledger]
\label{sum:pva-repair-ledger}
The repaired descent settles the foundational PVA layer as follows.

\begin{center}
\renewcommand{\arraystretch}{1.2}
\small
\begin{tabular}{@{}p{0.38\textwidth}|p{0.55\textwidth}@{}}
Issue & Repaired mechanism \\
\hline
Commutativity built into the definition & Use unsymmetrized $\mu$; prove commutativity from the exchange cylinder \\
Skew-symmetry only sketched & Use the same cylinder, with the degree-$(-1)$ transposition sign \\
Leibniz target unstable & Use the mixed three-face Stokes projection; recover the standard PVA rule \\
Jacobi mixed two-/three-face narratives & Use the full three-face singular Stokes argument \\
$\lambda$ vs.\ OPE variable conflated & Separate $\zeta$ and $\lambda$; connect them by $\mathcal B$
\end{tabular}
\end{center}

The higher contraction homotopies $K_k$ for $k\ge3$ are constructed via the affine retraction on $\R^{k-1}_{>0}$ in Proposition~4.22 of the main text.
\end{summary}
